% TOP_PROBLEM: {"id": 2303017, "title": "Research Problems in Function Theory — Problem 3.17", "category_id": 8, "category": {"id": 8, "name": "analysis", "display_name": "Analysis", "description": "Limits, continuity, calculus, and function theory.", "slug": "analysis", "order_index": 8, "created_at": "2026-07-31T15:26:25.671Z"}, "status": "open"}
% FIRST_POSTED: null
% ENTRY_KIND: statement_only
% Imported from: batch12.tex; UnsolvedMath ID 2303017
% Compile twice: lualatex 2303017.tex
\documentclass[11pt]{article}
\usepackage{fontspec}
\setmainfont{Latin Modern Roman}
\usepackage{amsmath,amssymb,mathtools,unicode-math}
\setmathfont{Latin Modern Math}
\usepackage{xurl,hyperref}
\hypersetup{colorlinks=true,urlcolor=blue,pdftitle={UnsolvedMath Problem 2303017}}
\newcommand{\sref}[2]{\hyperlink{source:#1:#2}{[S#2]}}
\newcommand{\cataloguescope}[1]{\paragraph{Mathematical question.}#1}
\begin{document}
% UnsolvedMath ID 2303017; source catalogue code AMR-022-3017
\setcounter{section}{2303016}
\section{Research Problems in Function Theory — Problem 3.17}\label{2303017}
{\small\sffamily UnsolvedMath ID 2303017 · Analysis}
\subsection{Definitions and mathematical statement}

\paragraph{Shared notation.}$\mathbb N=\{1,2,\ldots\}$, and $\mathbb Z,\mathbb Q,\mathbb R,\mathbb C$ denote the integers, rationals, reals and complex numbers. Any other conventions are specified locally.


A subharmonic function is an upper semicontinuous function with values in $[-\infty,\infty)$, not identically $-\infty$ on a connected component, that satisfies the mean-value inequality on every closed ball in its domain. A real harmonic function is a twice continuously differentiable function $u$ satisfying $\Delta u=0$. Let $D'\subset D$ be Lipschitz domains in $\mathbb R^n$, $n\geq3$: locally their boundaries are graphs of Lipschitz functions. Suppose $\partial D'\cap\partial D$ is a compact subset of the relative interior of $\Gamma\subset\partial D$. Fix $P_0\in D'$. Let $\mathcal H$ consist of positive harmonic functions on $D$ that extend continuously with value $0$ on $\Gamma$ and satisfy $h(P_0)=1$.
The source update attributes affirmative boundary-Harnack results to Dahlberg and Wu.
\paragraph{Statement recorded in the supplied source.}
Is there a finite constant $C$, depending only on $D,D',\Gamma,P_0$ and $n$, such that
\[h_1(P)\leq C h_2(P)\qquad(P\in D',\ h_1,h_2\in\mathcal H)?\] \sref{2303017}{2}
\subsection{Short English statement}
Are all normalized positive harmonic functions that vanish on the same boundary portion uniformly comparable throughout the smaller Lipschitz domain?
\subsection{Sources}
\begin{description}
\item[\hypertarget{source:2303017:1}{S1}] ULAM.AI and the UnsolvedMath Contributors, UnsolvedMath: A Curated Collection of Open Mathematics Problems, record 2303017 (AMR-022-3017). CC BY 4.0. \url{https://huggingface.co/datasets/ulamai/UnsolvedMath}.
\item[\hypertarget{source:2303017:2}{S2}] W. K. Hayman and E. F. Lingham, \emph{Research Problems in Function Theory}, arXiv:1809.07200v2 (2018), Problem 3.17 and Update 3.17. \url{https://arxiv.org/abs/1809.07200}.
\end{description}
\label{end:2303017}
\end{document}
